\section{Drei Zahlen, nicht eine}

\textbf{[SETZUNG]} Die Serie gibt zu jeder Vergleichszahl drei Angaben, nie eine:
den Wert, die Messgenauigkeit der Vergleichsgröße und die Bestimmungstoleranz
der Theoriegröße. Eine Prozentabweichung allein sagt nichts.

Der Grund ist einfach und wird durchgehend gebraucht: eine Abweichung von einem
Prozent ist bedeutungslos, wenn die Theoriegröße nur auf zwei Stellen bestimmt
ist -- und vernichtend, wenn beide Seiten auf acht Stellen feststehen. Ohne
beide Toleranzen ist die Prozentzahl nicht lesbar.

\section{Wir rechnen nichts aus}

\textbf{[SETZUNG]} Der wichtigste Satz über alle Zahlen dieser Serie: die Theorie
\emph{berechnet nichts}. Sie ist kein Rechenwerk, das Vorhersagen auf viele
Stellen produziert. Sie ordnet geometrische Strukturen und zeigt, dass sie in den
richtigen Bereich fallen.

Daraus folgt, wie die Prüfskripte zu lesen sind. Wenn ein Skript eine
Übereinstimmung \enquote{auf $10^{-15}$} meldet, ist das \emph{kein} Ergebnis. Es
ist entweder eine \textbf{Identität} -- eine Größe, die per Definition oder
Konstruktion gleich ist ($\alpha = 1$ nach der Ladungsnormierung, $L_0/\ell_P =
\xi$ nach der Definition) -- oder es ist eine \textbf{Bereichskontrolle}: die
Prüfung, ob wir in der richtigen Größenordnung liegen. Die vielen Nachkommastellen
sind Fließkomma-Arithmetik, nicht die Genauigkeit einer Aussage.

\textbf{[B]} \emph{Drei} Arten von Prüfung, streng zu trennen:

Die \emph{Identität} rechnet eine Definition nach. Sie kann gar nicht scheitern,
solange die Algebra stimmt. Sie zu bestehen ist kein Verdienst; sie zu verfehlen
wäre ein Programmierfehler. Beispiele: $\alpha = 4\pi/4\pi = 1$ (A130), $L_0 =
\xi\,\ell_P$ (A210), die Rückrichtung $\xi = 2\sqrt{Gm}$ (A145).

Die \emph{Bereichskontrolle} vergleicht eine Struktur mit der Messung und fragt
nur: richtige Größenordnung, richtiges Vorzeichen, richtige Skalierung? Ob
$E_\text{char}$ nahe $28$ statt bei $3$ oder $300$ liegt, ob $\kappa$ nahe $4$
statt bei $9$ liegt, ob der CHSH-Rest bei $10^{-4}$ statt $10^{-1}$ liegt.

Die \emph{Vorhersage} ist der Fall, in dem die Theorie \emph{sehr wohl} etwas
Messbares voraussagt: die \textbf{Leptonmassen}. Die Struktur $(\sqrt2, 2/9)$
legt $m_\tau$ aus $m_e$ und $m_\mu$ fest -- eine scharfe, falsifizierbare
Vorhersage, $m_\tau = 1776{,}968$ MeV gegen PDG $1776{,}86\pm0{,}12$, also
$+0{,}90\,\sigma$. Das ist keine Bereichskontrolle: hier steht ein Zahlenwert
\emph{vor} der Messung fest, und die Messung entscheidet ihn.

\textbf{[SETZUNG]} \enquote{Wir rechnen nichts aus} gilt also für die
\emph{Buchhaltungswerte} -- $\alpha = 1/137$, der SI-Wert von $G$, die Ankerwahl
$L_0$ --, nicht für die Leptonmassen. Bei den Buchhaltungswerten ist die Frage
\enquote{stimmt das Ergebnis exakt?} falsch gestellt; bei $m_\tau$ ist sie genau
richtig. Der Unterschied: die Leptonmassen sind durch eine \emph{Verhältnis}struktur
so eng gebunden, dass die dritte Masse aus zweien folgt -- das ist eine echte
Vorhersage, kein Absolutwert und keine Kalibrierung.

\section{Die vier Genauigkeitsklassen der Serie}

\textbf{[B]} Über die ganze Serie treten vier Klassen auf. Sie zu unterscheiden
ist die eigentliche Arbeit des Toleranzmaßstabs.

\textbf{Klasse 1 -- fünf Stellen und mehr.} Der Zirkulant (A110): $Q=2/3$ auf
sechs Stellen, $\theta(\mu/e)$ auf sieben. Hier ist eine Abweichung von
$10^{-6}$ ein ernstes Signal, weil beide Seiten so genau feststehen.

\textbf{Klasse 2 -- Prozentniveau.} Die Leiter (A100): $\mu/e$ auf $0{,}52\,\%$,
$\tau/\mu$ auf $1{,}04\,\%$. Hier ist ein Prozent \emph{kein} Signal gegen die
Struktur, weil die Struktur auf diesem Niveau nichts Genaueres behauptet -- und
kein Erfolg, den man auf fünf Stellen ausweisen dürfte.

\textbf{Klasse 3 -- Buchhaltungsgenauigkeit.} Der SI-Wert von $\alpha$ (A130,
$-0{,}24\,\%$) und von $G$ (A145, $E_\text{char}$ auf $1{,}5\,\%$). Diese
Abweichungen betreffen die Skalenumrechnung, nicht die Struktur. Sie sind kein
Test (A210); die Prozentzahl misst die Brücke, nicht die Theorie.

\textbf{Klasse 4 -- unbestimmt.} Der Neutrinosektor (A150), die Quarkmassen
(A150, schemaabhängig). Hier ist die Theoriegröße so wenig festgelegt, dass jede
Übereinstimmung folgenlos ist. Was nicht widerlegt werden kann, bestätigt nichts.

\section{Die häufigsten Fehlbuchungen}

\textbf{[B]} Der Toleranzmaßstab existiert, um fünf Fehler zu vermeiden, die
alle im Korpus mindestens einmal vorkamen:

\textbf{Klasse verwechseln.} Eine Klasse-2-Abweichung (Leiter, Prozent) gegen
eine Klasse-1-Genauigkeit (Zirkulant, fünf Stellen) halten -- das ist der
Ebenenfehler aus A080/A110. Beide beschreiben dieselben Massen, aber auf
verschiedenen Niveaus.

\textbf{Buchhaltung als Test lesen.} Eine Klasse-3-Abweichung als Bestätigung
oder Widerlegung zählen. Der Fehler, den A130 an der $\alpha$-Darstellung
benennt.

\textbf{Unbestimmtes als bestätigt zählen.} Eine Klasse-4-Größe als Erfolg
buchen, weil die Theorie \enquote{im Bereich} liegt -- obwohl der Bereich so
weit ist, dass alles hineinfällt.

\textbf{Stellen erfinden.} Einer Klasse-2-Größe fünf Stellen geben, weil die
Rechnung sie liefert. Die Fließkommagenauigkeit einer Rechnung ist nicht die
Bestimmungstoleranz der Größe.

\textbf{Abhängiges als unabhängig zählen.} Eine Größe als eigene Evidenz oder als
zweite Kennzahl führen, obwohl sie aus einer bereits gezählten \emph{arithmetisch
folgt}. Beispiel aus der Leiter (A100): im Verhältnis
$N_1{:}N_2{:}N_3 \approx 1{:}1{,}981{:}2{,}971$ ist allein die $1{,}981$ gemessen
-- die dritte Zahl ist über $p_3 = p_1 + p_2$ (bei durchgängig multiplikativer
Korrektur $K^{p_3} = K^{p_1}K^{p_2}$) als $1 + 1{,}981$ erzwungen und trägt keine
unabhängige Evidenz. Ebenso ist die Toleranz-Bilanzzahl \enquote{Faktor rund
$300$} keine zweite Größe, sondern der Kehrwert derselben Abweichung ($N_2/N_1$
von exakt $2$). Das generationslineare Gesetz $N_g = g\,N_0$ ruht damit auf
\emph{einer} Beobachtung, nicht auf dreien; es ist ein Kandidat auf einer
Messung, keine dreifach gestützte Regularität. Der Maßstab
verlangt, solche erzwungenen Stellen als das zu buchen, was sie sind -- Folge,
nicht Fund.

\section{Der Maßstab in einem Satz}

\textbf{[SETZUNG]} Eine Übereinstimmung zählt nur so weit, wie die
\emph{schlechter} bestimmte der beiden Seiten reicht. Eine Abweichung wiegt nur
so schwer, wie die \emph{besser} bestimmte der beiden Seiten sie trägt.

Das ist der ganze Maßstab. Alles andere ist seine Anwendung.

\section{Prüfskript}

\begin{itemize}
  \item Einordnung der Vergleichszahlen der Serie in die vier Klassen, mit
    Ausweis von Messgenauigkeit und Bestimmungstoleranz je Fall; und die
    Gegenprobe, dass eine Klasse-2-Zahl auf fünf Stellen ausgewiesen eine
    Scheingenauigkeit erzeugt:
    \texttt{python/A\_Serie\_Skripte/a220\_toleranzmassstab.py}
\end{itemize}

\section{Was offen bleibt}

\textbf{Erstens ist die Bestimmungstoleranz mancher Theoriegrößen selbst nicht
scharf.} Für die Leiterexponenten (A100) ist unklar, auf wie viele Stellen sie
überhaupt bestimmt sind. Wo das so ist, ist der Maßstab konservativ anzuwenden.

\textbf{Zweitens ersetzt der Maßstab keine Fehlerrechnung.} Er ordnet
Größenordnungen, er propagiert keine Unsicherheiten durch Formeln. Eine volle
Fehlerfortpflanzung leistet die Serie nicht.

\section{Ersetzt}

Dieses Dokument macht die Toleranzregel aus A010 und A080 zum eigenen Prinzip
und wendet sie auf alle Sektoren an. Es nimmt die entsprechenden Teile auf von
Dok.~122 (Verhältnis und Absolutes) und Dok.~292 (Toleranzanalyse Leptonen).
Eingearbeitet: R55 (Evidenzbreite der Leiterkopplung: eine Beobachtung statt
drei; \enquote{Faktor 300} nicht als Kennzahl).

\section{Verwendung von KI-Werkzeugen}

Dieses Dokument ist unter Verwendung KI-gestützter Werkzeuge entstanden. Die
Fragestellungen, die Setzungen und alle inhaltlichen Entscheidungen stammen vom
Autor; die Werkzeuge formulieren aus, rechnen nach und erstellen Prüfskripte.
Die Verantwortung für Inhalt und Fehler liegt vollständig beim Autor. Der
ausführliche Wortlaut steht in A010.
